% ==================================================================
%  Relative validity criteria --- MaxCriterion and ElbowCriterion.
%
%  Copied from the measure template
%  (documentation/sections/clustering_methods/00-template.tex is the
%  method template; the measure template is the one measure_1.tex was
%  made from). A relative criterion is a validity index, so the measure
%  parts apply; the index-specific prompts are kept and the
%  dissimilarity prompts dropped.
%
%  Add one \input line for this file under Section 7.1 in
%  cluster_main.tex, beside measure_1.tex.
%
%  Code counterpart: xxcluster/measures/validation/relative.py.
%  Part 3 (Properties) agrees with the class attributes declared there.
%
%  NOTE: citation keys below are written as ref_<n> placeholders. Before
%  merge, replace each with the real key from the shared citation-key
%  mapping sheet and record it in the REFERENCES USED block. The keys
%  here are marked TODO so they are not mistaken for real ones.
% ==================================================================
\subsubsection{Relative criteria: taking a choice from a curve}
\label{sec:measure:relative}

\paragraph{Overview}
\label{sec:measure:relative:overview}
A relative criterion is a validity index in the third sense of
Section~\ref{sec:clustering_methods}: it does not score one partition but
reads a sequence of scores and selects from it. The sequence is produced
by evaluating an internal index over a family of candidate partitions,
typically the same method at each value of the cluster count~$k$, and the
criterion returns the~$k$ the curve points to. This is the group the
choice of~$k$ rests on: the peak of a silhouette curve, or the elbow of a
sum-of-squares curve. Two criteria are implemented. \textsc{MaxCriterion}
takes the best-scoring candidate; \textsc{ElbowCriterion} takes the
candidate of greatest curvature, for the monotone curves on which a
maximum is uninformative. The selection rule is deliberately kept separate
from the index it reads, since the same curve yields different answers
under different rules, and which rule was used is part of the reported
result.

\paragraph{Definition}
\label{sec:measure:relative:definition}
Let $\mathcal{K}$ be the swept candidates and $s\colon \mathcal{K}\to
\mathbb{R}$ the score of each under a fixed internal index, with the
index's direction given by a relation $\succ$ (so $a \succ b$ reads ``$a$
is better than $b$''). \textsc{MaxCriterion} returns
$\arg\max_{k\in\mathcal{K}}^{\succ} s(k)$, the optimum \emph{in the
direction the index declares} rather than the numerical maximum.
\textsc{ElbowCriterion} is defined geometrically. Order the candidates,
form the points $(k, s(k))$, and rescale both axes to the unit square.
Let $\ell$ be the chord from the first point to the last --- the curve a
constant marginal return would trace. The elbow is
$\arg\max_{k\in\mathcal{K}} \operatorname{dist}\big((k, s(k)),\, \ell\big)$,
the candidate whose normalised score lies furthest from the chord. The
knee/elbow construction follows the distance-to-chord method of
\cite{ref_TODO_satopaa_kneedle}; the use of a criterion curve to choose~$k$
is standard, e.g.\ \cite{ref_TODO_thorndike_elbow}.

\paragraph{Properties}
\label{sec:measure:relative:properties}
These agree with the class attributes in \texttt{relative.py}.
\textsc{MaxCriterion} carries no direction of its own: it selects through
the index's \texttt{is\_better}, so a maximised index (silhouette) yields
the peak and a minimised index (Davies--Bouldin) the trough, with no code
change. It requires the index at construction, with no default, for the
same reason an index requires its direction to be declared --- a guessed
direction eventually selects the worst candidate and reports it as the
best. Ties break by sweep order, and since \texttt{is\_better} is strict
this returns the smallest~$k$ achieving the best score, the conservative
choice. \textsc{ElbowCriterion} is direction-free: distance from a chord
does not depend on whether the index is maximised or minimised, which is
why an elbow can be read from a minimised curve without inverting it. Both
axes are normalised before the distance is taken, so the selection is
invariant to the index's units and to the width of the sweep. A failed run
enters the curve as \texttt{NaN}; it is retained so the run is still
reported, and \texttt{is\_better} guarantees a \texttt{NaN} never wins.

\paragraph{Applicability}
\label{sec:measure:relative:applicability}
A relative criterion applies wherever a family of partitions is compared
on a common index. \textsc{MaxCriterion} accepts any candidate type, since
it only compares scores. \textsc{ElbowCriterion} requires numeric
candidates, because it measures a distance along the swept axis; a sweep
over a non-numeric parameter raises rather than returning a meaningless
answer. Both consume a soft partition only once it has been defuzzified to
the hard labels their underlying index scores. Noise handling is inherited
from that index: where the index excludes noise-labelled observations, so
does the curve, and the exclusion must be stated wherever the selection is
reported.

\paragraph{Computation and complexity}
\label{sec:measure:relative:complexity}
Given the curve, both criteria are linear in the number of candidates
$|\mathcal{K}|$: \textsc{MaxCriterion} is a single pass, and
\textsc{ElbowCriterion} a sort followed by one pass of point-to-line
distances. The cost of the procedure is therefore dominated not by the
criterion but by producing the curve --- one clustering fit and one index
evaluation per candidate. Where the index requires the full pairwise
distance matrix (silhouette does), that evaluation is quadratic in the
number of observations, and at \projectname{} data volumes it is the
binding constraint, not the selection.

\paragraph{Application to \projectname{} data}
\label{sec:measure:relative:application}
The criteria are implemented in
\texttt{xxcluster/measures/validation/relative.py} and registered as
\texttt{"max"} and \texttt{"elbow"}. They decide only; the procedure that
generates the candidate partitions and applies a criterion is
\texttt{xxcluster.selection.n\_clusters}, and the curve a criterion
returns is drawn by \texttt{plot\_selection\_curve} in
\texttt{xxcluster.viz.diagnostics}, which marks the selection and draws the
curve whether or not it was conclusive. \textsc{MaxCriterion} carries two
declared thresholds, \texttt{flat\_tolerance} and \texttt{min\_margin},
each a fraction of the index's declared range (for silhouette, the span
from $-1$ to $1$); \textsc{ElbowCriterion} carries \texttt{min\_curvature},
a fraction of the chord length. They are parameters rather than constants
so that the write-up can state the threshold a selection had to clear. The
accompanying notebook evidences the behaviour end to end (Appendix, notebook
\texttt{task37\_relative\_criteria.ipynb}).

\paragraph{Behaviour}
\label{sec:measure:relative:behaviour}
Demonstrated on the \texttt{iris} benchmark rather than asserted, since
its structure is known. Sweeping $k=2,\dots,10$, the silhouette curve peaks
at $k=2$ and \textsc{MaxCriterion} selects it; replacing the index with a
minimised copy of itself makes the identical criterion select the trough at
$k=9$, confirming the direction is the index's. The within-cluster
sum-of-squares (inertia) curve falls monotonically, so its numerical best
is always the largest~$k$ swept; \textsc{MaxCriterion} on it returns the
sweep boundary $k=10$, while \textsc{ElbowCriterion} returns the bend at
$k=4$. This is the behaviour the pair exists to provide: a maximum where
the curve has a peak, and an elbow where it does not.

\paragraph{Limitations}
\label{sec:measure:relative:limitations}
A criterion can always return an answer, but the answer is not always
warranted, and \texttt{is\_conclusive} reports when it is not. For
\textsc{MaxCriterion} the curve is inconclusive when there are fewer than
three candidates, when the spread is below \texttt{flat\_tolerance} (the
candidates are indistinguishable), when the winner beats the runner-up by
less than \texttt{min\_margin} (a tie, not a peak), or when the curve is
monotone with its best value at a sweep boundary (an artefact of where the
sweep stopped, to be answered by widening it, which is exactly the case the
elbow handles). For \textsc{ElbowCriterion} the curve is inconclusive when
its greatest normalised departure from the chord is below
\texttt{min\_curvature} --- a straight line has no elbow, and reporting one
turns rounding into a finding. A further limitation is inherited: the
criterion is only as neutral as the index beneath it, and a shape-biased
index (Section~\ref{sec:measure:relative:behaviour}) carries that bias into
the selection, a threat to validity for Section~\ref{sec:methodology:threats}
rather than a detail.

\paragraph{Summary}
\label{sec:measure:relative:summary}
\textsc{MaxCriterion} is the default rule for a peaked curve and
\textsc{ElbowCriterion} the rule for a monotone one; both defer the
direction and the reporting to the index they read, and both declare when
the curve does not support a choice. They are part of the reported protocol
of Section~\ref{sec:methodology:metrics}, being the mechanism by which~$k$
is chosen wherever a sweep is reported.

% ==================================================================
%  REFERENCES USED --- fill this in before merge.
%
%  Replace each ref_TODO_* above with the real ref_<n> key from the
%  shared citation-key mapping sheet, then list them here, one per line,
%  including anything read but not cited. The same keys go in the
%  `references' field where MaxCriterion/ElbowCriterion are registered
%  in xxcluster (CONTRIBUTING.md, 2.3.1); the two lists must agree.
%
%    ref_??   Thorndike (1953), "Who belongs in the family?"   cited in: Definition
%    ref_??   Satopää et al. (2011), Kneedle / knee detection  cited in: Definition
% ==================================================================
